% ----- CHAUPAI 8 -----

\newpage

\subsection*{\deva{अष्टम चौपाई} (Eighth Chaupai)}
\addcontentsline{toc}{subsection}{\deva{अष्टम चौपाई} (Eighth Chaupai) — \deva{प्रभु चरित्र...}}

\begin{minipage}[t]{0.55\textwidth}
\vspace{0pt}

{\large \linespread{1.5}\selectfont
\deva{प्रभु चरित्र सुनिबे को रसिया।}\\
\deva{राम लखन सीता मन बसिया॥}
\par}

\vspace{0.5em}

{\large \linespread{1.5}\selectfont
\telu{ప్రభు చరిత్ర సునిబే కో రసియా।}\\
\telu{రామ లఖన సీతా మన బసియా॥}
\par}

\vspace{0.5em}

{\large \linespread{1.3}\selectfont
\textit{prabhu caritra sunibe ko rasiyā |}\\
\textit{rāma lakhana sītā mana basiyā ||}
\par}
\end{minipage}%
\hfill
\begin{minipage}[t]{0.42\textwidth}
\vspace{0pt}
\begin{tcolorbox}[anchorbox]
\textbf{The Art of Listening:} He is praised for being exceptionally eager to listen (\textit{Sunibe ko Rasiya}). In the *Valmiki Ramayana* (*Kishkindha Kanda, 3:28-33*), Lord Rama notes that Hanuman's flawless speech is proof of his highly disciplined, attentive listening as a student of the Vedas. True mastery in communication and leadership begins with the patient, active absorption of knowledge.
\vspace{0.5em}\newline Hanuman's legendary eloquence was the direct result of his deep, attentive listening and study under his mentors.\newline\null\hfill\href{https://www.valmiki.iitk.ac.in/sloka?field_kanda_tid=4&language=dv&field_sarga_value=3}{\textit{Kishkindha Kanda, 3:28-33}}
\end{tcolorbox}
\end{minipage}

\vspace{0.5em}
\subsubsection*{\textbf{Visual Rhythm Map:}}
\begin{table}[h!]
\centering
\renewcommand{\arraystretch}{1.2}
\setlength{\tabcolsep}{0pt}
\begin{tabularx}{\textwidth}{|*{16}{>{\centering\arraybackslash\hsize=1.0\hsize}X|}}
\hline
\multicolumn{2}{|>{\centering\arraybackslash\hsize=2.0\hsize}X|}{\textbf{\laghu{प्र}\laghu{भु}} \par (2)} &
\multicolumn{4}{>{\centering\arraybackslash\hsize=4.0\hsize}X|}{\textbf{\laghu{च}\guru{रि}\laghu{त्र}} \par (4)} &
\multicolumn{4}{>{\centering\arraybackslash\hsize=4.0\hsize}X|}{\textbf{\laghu{सु}\laghu{नि}\guru{बे}} \par (4)} &
\multicolumn{2}{>{\centering\arraybackslash\hsize=2.0\hsize}X|}{\textbf{\guru{को}} \par (2)} &
\multicolumn{4}{>{\centering\arraybackslash\hsize=4.0\hsize}X|}{\textbf{\laghu{र}\laghu{सि}\guru{या}} \par (4)} \\
\hline
\end{tabularx}
\vspace{-1.5em}
\end{table}

\begin{table}[h!]
\centering
\renewcommand{\arraystretch}{1.2}
\setlength{\tabcolsep}{0pt}
\begin{tabularx}{\textwidth}{|*{16}{>{\centering\arraybackslash\hsize=1.0\hsize}X|}}
\hline
\multicolumn{3}{|>{\centering\arraybackslash\hsize=3.0\hsize}X|}{\textbf{\guru{रा}\laghu{म}} \par (3)} &
\multicolumn{3}{>{\centering\arraybackslash\hsize=3.0\hsize}X|}{\textbf{\laghu{ल}\laghu{ख}\laghu{न}} \par (3)} &
\multicolumn{4}{>{\centering\arraybackslash\hsize=4.0\hsize}X|}{\textbf{\guru{सी}\guru{ता}} \par (4)} &
\multicolumn{2}{>{\centering\arraybackslash\hsize=2.0\hsize}X|}{\textbf{\laghu{म}\laghu{न}} \par (2)} &
\multicolumn{4}{>{\centering\arraybackslash\hsize=4.0\hsize}X|}{\textbf{\laghu{ब}\laghu{सि}\guru{या}} \par (4)} \\
\hline
\end{tabularx}
\vspace{-1.5em}
\end{table}

\vspace{1em}
\textbf{ \deva{सरल-संस्कृतम्}:}\\
 \deva{भवान् प्रभोः (श्रीरामस्य) चरित्रं श्रोतुं रसिकः अस्ति।}\\
 \deva{श्रीरामः, लक्ष्मणः, सीता च भवतः मनसि वसन्ति।}\\


You delight in listening to the stories and glories of the Lord.\\
Rama, Lakshmana, and Sita reside in your mind.


\vspace{1em}
\newpage
\textbf{\deva{प्रतिपदार्थः} (Meaning of each word):}
\begin{table}[h!]
\centering
\renewcommand{\arraystretch}{1.2}
\begin{tabularx}{\textwidth}{@{} l l X X @{}}
\toprule
\textbf{Awadhi} & \textbf{Sanskrit} & \textbf{English} & \textbf{Grammar \& Notes} \\
\midrule
\deva{प्रभु} / \telu{ప్రభు} / \textit{prabhu}  &  \deva{प्रभोः}  &  Lord's  &   \\
\deva{चरित्र} / \telu{చరిత్ర} / \textit{caritra}  &  \deva{चरित्रम् / कथाम्}  &  Character / Stories  &   \\
\deva{सुनिबे को} / \telu{సునిబే కో} / \textit{sunibe ko}  &  \deva{श्रोतुम्}  &  To listen  & Awadhi infinitive verb \\
\deva{रसिया} / \telu{రసియా} / \textit{rasiyā}  &  \deva{रसिकः}  &  One who relishes  &   \\
\deva{राम} / \telu{రామ} / \textit{rāma}  &  \deva{श्रीरामः}  &  Rama  &   \\
\deva{लखन} \gotchaicon / \telu{లఖన} / \textit{lakhana}  &  \deva{लक्ष्मणः}  &  Lakshmana  & cluster simplification \\ 
\deva{सीता} \gotchaicon / \telu{సీతా} / \textit{sītā}  &  \deva{सीता} \gotchaicon  &  Sita  &   \\
\deva{मन} / \telu{మన} / \textit{mana}  &  \deva{मनसि}  &  In the mind  &   \\
\deva{बसिया} \gotchaicon / \telu{బసియా} / \textit{basiyā}  &  \deva{వसन्ति / వాసీ}  &  Reside / Dweller  & \\ 
\bottomrule
\end{tabularx}
\end{table}

\begin{tcolorbox}[gotchabox]
\begin{itemize}
    \item \textbf{Consonant Melt (\deva{लक्ष्मण $\rightarrow$ लखन}):} The harsh consonant cluster `kshma` (\deva{क्ष्म}) in Sanskrit melts down into the softer, breathy `kha` (\deva{ख}) in Awadhi, making it much easier to sing! The Gita Press version uses \deva{लषन} instead of \deva{लखन}. Both carry the same metrical value.
    \item \textbf{The \deva{व} $\rightarrow$ \deva{ब} Shift (\deva{बसिया}):} The Sanskrit root for dwelling is \deva{वस्} (Vas). True to form, the Awadhi immediately hardens that initial $v$ into a $b$, giving us \deva{बसिया} (Basiyaa).
    
\end{itemize}
\end{tcolorbox}



\newpage
